Se dice que $f: \mathbb{R}^2 \to \mathbb{R}$ es homogénea de grado $n$ si $f(tx, ty) = t^n f(x, y)$ para todo $t > 0$.
\begin{enumerate}
    \item Demuestre el teorema de Euler para funciones homogéneas: si $f$ es diferenciable y homogénea de grado $n$, entonces $x f_x + y f_y = n f(x, y)$. Sugerencia: diferencie la identidad $f(tx, ty) = t^n f(x,y)$ con respecto a $t$ y evalúe en $t = 1$.
    \item Verifique el teorema para $f(x, y) = x^3 + 2x^2y - xy^2$.
    \item Demuestre que si $f$ es homogénea de grado $n$ y tiene derivadas parciales de segundo orden continuas, entonces $x^2 f_{xx} + 2xy f_{xy} + y^2 f_{yy} = n(n-1)f$.
\end{enumerate}
